\chapter{FPGA 仿真与验证}
\label{chap:control}
\label{chap:fpga-simulation}

本章围绕“输入数据是否正确、帧配置是否一致、访存是否遵守接口约定、结果是否完整提交”
建立 RTL 仿真证据链。验证对象包括 IMU 接收、姿态跨域、像素预处理、运行时控制、
逆映射与分块渲染，以及共享 DDR 访问。2026 年 10 月 25 日，使用工程现有测试平台
复跑了 11 项功能测试，全部通过；其中分块渲染采用缩小的输出区域，以覆盖块边界和异常处理。
完整 $1920\times1080$ 并发链路和厂商 DDR 模型结果分别引用已有归档，明确保留其版本与条件。
本章给出功能验证方法与观测结果，实时性能和资源效率见第~\ref{chap:verification}~章，
矫正清晰度、直线度及边缘质量的实测方案见第~\ref{chap:performance}~章。

\section{仿真环境与验证层次}
\subsection{工具、时钟与证据管理}
本次复跑使用 Icarus Verilog 12.0，以 SystemVerilog 模式编译测试平台，
由 \texttt{vvp} 执行。复跑脚本直接读取 \path{../Pgl50h/tests/} 与
\path{../Pgl50h/src/} 中的源码，不改写 FPGA 工程；编译产物存放于临时目录，
日志、参数、输入文件 SHA256 及选定信号的 VCD 保存在本报告的
\path{chapter5/references/local_20261025/}。清单逐项记录编译及运行命令，
因此本次结果对应清单中的源码内容，而非仅对应一个工程目录名。

测试平台按模块使用不同的时钟：IMU 和控制域通常为 50 MHz，
视频像素域以半周期 3.367 ns 近似 148.5 MHz，DDR 用户侧及主要渲染模块为 100 MHz。
各平台分别检查复位、有效信号和跨域事件，不将独立时钟视为同相时钟。
\path{tb_pc_control} 为缩短串口测试时间，将模块时钟频率参数设为 1.8432 MHz，
保持 115200 baud 的整数分频关系；其结果用于验证协议状态机，
不能替代真实 50 MHz 时钟和板级串口电气测试。

\begin{table}[htbp]\centering\small
\caption{本章采用的三层仿真证据}\label{tab:sim-evidence}
\begin{tabularx}{\textwidth}{>{\raggedright\arraybackslash}p{2.5cm}YY}
\toprule
层次与日期 & 模型及覆盖范围 & 结果的适用边界 \\
\midrule
本次功能复跑\newline 10 月 25 日 & 当前读取的 RTL、原工程测试平台；11 项测试与两组 VCD & 证明所选激励下的功能和协议行为；未重跑完整 1080p 并发链路 \\
整帧并发归档\newline 10 月 22 日 & ModelSim、采集/渲染/显示并发、共享 DDR 行为模型 & 对应归档源码快照；吞吐结果不能自动转移到后续修改版本 \\
厂商 DDR 归档\newline 10 月 10 日 & 实际生成的控制器/PHY 与两颗 Micron x16 模型 & 验证短事务的数据、边界及返回时序；不代表整帧并发或实板结果 \\
\bottomrule
\end{tabularx}
\end{table}

\subsection{自动判定方法与测试矩阵}
各测试平台通过断言和计分板自动判定结果：输入激励产生独立期望值或事务序列，
输出有效时逐项比较；超时、数据错误、重复提交和非法协议均触发失败。
复跑脚本要求编译和运行返回码均为零、日志包含 \texttt{PASS} 且无
\texttt{FATAL/ERROR}。本次 11 项测试的编译日志均为空，运行均满足上述条件。
其中 DDR 仲裁和渲染测试主动注入非法响应，预期行为是模块报告错误并阻止不安全完成，
测试平台检查该行为后仍以 \texttt{PASS} 结束。

\begin{table}[htbp]\centering\small
\caption{本次实际执行的功能测试矩阵（全部通过）}\label{tab:sim-matrix}
\begin{tabularx}{\textwidth}{>{\raggedright\arraybackslash}p{4.5cm}Y}
\toprule
测试平台（省略 \texttt{tb\_} 前缀） & 主要检查内容与规模 \\
\midrule
\path{imu_protocol} & 921600 baud，符号与字节序、坏校验拒收、重复帧头重同步 \\
\path{attitude_cdc} & 100 组异步姿态快照，无混包、无重复，目的域复位 \\
\path{pose_async_fifo} & 512/8/2 深度、独立时钟、满空及回绕、反压与非空复位 \\
\path{pose_ddr_arbiter} & \texttt{STREAM=1}；160 笔写/760 beats、128 笔读/488 beats \\
\path{pose_inverse_project_q14} & 默认除法路径；姿态就绪、坐标参考点、限幅和无效映射 \\
\path{adaptive_gamma_stream} & 2048 个 LUT 表项，15365 个有效像素、60 次帧统计 \\
\path{pc_control} & 8 种掩码、SOF 应用、busy、校验、超时、查询和遥测 \\
\path{gamma_telemetry} & \texttt{COLOR=1}；彩色 Gamma 参考、前后亮度和、跨域快照 \\
\path{pc_control_frame_alignment} & \texttt{GAMMA=1}；25 帧、8 种掩码、SOF 与首像素对齐 \\
\path{rgb888_raster} & 全部 65536 种 RGB565 输入，DE 转发，23/2 拍光栅延迟与复位 \\
\path{pose_renderer_tiled} & $144\times40$ 输出区域，$128\times32$ 块，5 组姿态与异常 LAST \\
\bottomrule
\end{tabularx}
\end{table}
\sourcefile{\path{chapter5/references/run_simulations.py}；本次日志与源码哈希见
\path{chapter5/references/local_20261025/manifest.json}}

\section{姿态输入与跨时钟域验证}
\subsection{IMU 串口帧解析}
IMU 测试直接向接收端串行注入 8-N-1 数据。帧为
\begin{center}
\texttt{55 AA | yaw[2] | pitch[2] | roll[2] | SUM | 0D}，
\end{center}
其中角度为 16 bit 小端有符号 Q2048 弧度，校验为六个载荷字节的模 256 加和。
测试使用 50 MHz 接收时钟与约 1085 ns 位周期，依次发送合法帧、校验错误帧、
噪声字节及包含重复 \texttt{55} 的重新同步序列。

第一组合法载荷为 yaw=\texttt{FC00}、pitch=\texttt{0200}、roll=\texttt{FF00}，
检查负数不被解释为无符号大角度，且三个字段不会发生大小端交换。
坏校验帧之后，姿态寄存器和事件翻转标志均应保持不变。
噪声与重复帧头之后，再发送 yaw=\texttt{0100}、pitch=\texttt{FE00}、
roll=\texttt{0080} 的合法帧，检查接收器能够恢复。
测试最终只允许两次合法更新，表明异常字节没有形成虚假的姿态事件。

\subsection{异步姿态快照与波形分析}
\path{attitude_cdc} 的源域为 50 MHz，目的域约为 148.5 MHz。
测试生成 100 组不断变化的 yaw、pitch、roll，源端以翻转标志通知更新，
目的端每次 \path{valid} 有效时同时比较三个字段，并统计接收次数。
因此，单独检查某一个角度正确还不够，三个字段必须来自同一组输入快照。
结束前复位目的域，检查输出数据和有效信号被清零。

图~\ref{fig:sim-cdc} 直接由本次 VCD 绘制，显示第一组快照：
160 ns 时源端更新 yaw=\texttt{000D}、pitch=\texttt{FFF2}、
roll=\texttt{B9B0}，并翻转事件标志；
178.451 ns 时目的端三个输出同时更新，\path{valid} 保持一个目的时钟周期，
于 185.185 ns 撤销。该片段的事件到输出延迟为 18.451 ns，
它只对应本次时钟相位，不能作为所有异步相位下固定的跨域延迟。

\begin{figure}[htbp]\centering
% Generated from actual VCD by references/make_waveforms.py.
\begin{tikzpicture}[x=1cm,y=1cm,font=\scriptsize]
\tikzset{trace/.style={draw=TemplateBlue,line width=0.65pt}}
\draw[gray!25] (0.00000,-5.550) -- (0.00000,0.75);
\node[anchor=south] at (0.00000,0.76) {140};
\draw[gray!25] (1.57143,-5.550) -- (1.57143,0.75);
\node[anchor=south] at (1.57143,0.76) {150};
\draw[gray!25] (3.14286,-5.550) -- (3.14286,0.75);
\node[anchor=south] at (3.14286,0.76) {160};
\draw[gray!25] (4.71429,-5.550) -- (4.71429,0.75);
\node[anchor=south] at (4.71429,0.76) {170};
\draw[gray!25] (6.28571,-5.550) -- (6.28571,0.75);
\node[anchor=south] at (6.28571,0.76) {180};
\draw[gray!25] (7.85714,-5.550) -- (7.85714,0.75);
\node[anchor=south] at (7.85714,0.76) {190};
\draw[gray!25] (9.42857,-5.550) -- (9.42857,0.75);
\node[anchor=south] at (9.42857,0.76) {200};
\draw[gray!25] (11.00000,-5.550) -- (11.00000,0.75);
\node[anchor=south] at (11.00000,0.76) {210};
\node[anchor=west] at (11.15,0.87) {ns};
\node[anchor=east] at (-0.15,0.160) {\texttt{src\_clk}};
\draw[trace] (0.00000,0.000) -- (1.57143,0.000) -- (1.57143,0.320) -- (3.14286,0.320) -- (3.14286,0.000) -- (4.71429,0.000) -- (4.71429,0.320) -- (6.28571,0.320) -- (6.28571,0.000) -- (7.85714,0.000) -- (7.85714,0.320) -- (9.42857,0.320) -- (9.42857,0.000) -- (11.00000,0.000);
\node[anchor=east] at (-0.15,-0.440) {\texttt{dst\_clk}};
\draw[trace] (0.00000,-0.280) -- (0.22220,-0.280) -- (0.22220,-0.600) -- (0.75130,-0.600) -- (0.75130,-0.280) -- (1.28040,-0.280) -- (1.28040,-0.600) -- (1.80950,-0.600) -- (1.80950,-0.280) -- (2.33860,-0.280) -- (2.33860,-0.600) -- (2.86770,-0.600) -- (2.86770,-0.280) -- (3.39680,-0.280) -- (3.39680,-0.600) -- (3.92590,-0.600) -- (3.92590,-0.280) -- (4.45500,-0.280) -- (4.45500,-0.600) -- (4.98410,-0.600) -- (4.98410,-0.280) -- (5.51320,-0.280) -- (5.51320,-0.600) -- (6.04230,-0.600) -- (6.04230,-0.280) -- (6.57140,-0.280) -- (6.57140,-0.600) -- (7.10050,-0.600) -- (7.10050,-0.280) -- (7.62960,-0.280) -- (7.62960,-0.600) -- (8.15870,-0.600) -- (8.15870,-0.280) -- (8.68780,-0.280) -- (8.68780,-0.600) -- (9.21690,-0.600) -- (9.21690,-0.280) -- (9.74600,-0.280) -- (9.74600,-0.600) -- (10.27510,-0.600) -- (10.27510,-0.280) -- (10.80420,-0.280) -- (10.80420,-0.600) -- (11.00000,-0.600);
\node[anchor=east] at (-0.15,-1.040) {\texttt{toggle}};
\draw[trace] (0.00000,-1.200) -- (3.14286,-1.200) -- (3.14286,-0.880) -- (11.00000,-0.880);
\node[anchor=east] at (-0.15,-1.640) {\texttt{sy}};
\draw[trace] (0.00000,-1.640) -- (0.06000,-1.480) -- (3.08286,-1.480) -- (3.14286,-1.640) -- (3.08286,-1.800) -- (0.06000,-1.800) -- cycle;
\node at (1.57143,-1.640) {\texttt{0000}};
\draw[trace] (3.14286,-1.640) -- (3.20286,-1.480) -- (10.94000,-1.480) -- (11.00000,-1.640) -- (10.94000,-1.800) -- (3.20286,-1.800) -- cycle;
\node at (7.07143,-1.640) {\texttt{000D}};
\node[anchor=east] at (-0.15,-2.240) {\texttt{dy}};
\draw[trace] (0.00000,-2.240) -- (0.06000,-2.080) -- (5.98230,-2.080) -- (6.04230,-2.240) -- (5.98230,-2.400) -- (0.06000,-2.400) -- cycle;
\node at (3.02115,-2.240) {\texttt{0000}};
\draw[trace] (6.04230,-2.240) -- (6.10230,-2.080) -- (10.94000,-2.080) -- (11.00000,-2.240) -- (10.94000,-2.400) -- (6.10230,-2.400) -- cycle;
\node at (8.52115,-2.240) {\texttt{000D}};
\node[anchor=east] at (-0.15,-2.840) {\texttt{sp}};
\draw[trace] (0.00000,-2.840) -- (0.06000,-2.680) -- (3.08286,-2.680) -- (3.14286,-2.840) -- (3.08286,-3.000) -- (0.06000,-3.000) -- cycle;
\node at (1.57143,-2.840) {\texttt{0000}};
\draw[trace] (3.14286,-2.840) -- (3.20286,-2.680) -- (10.94000,-2.680) -- (11.00000,-2.840) -- (10.94000,-3.000) -- (3.20286,-3.000) -- cycle;
\node at (7.07143,-2.840) {\texttt{FFF2}};
\node[anchor=east] at (-0.15,-3.440) {\texttt{dp}};
\draw[trace] (0.00000,-3.440) -- (0.06000,-3.280) -- (5.98230,-3.280) -- (6.04230,-3.440) -- (5.98230,-3.600) -- (0.06000,-3.600) -- cycle;
\node at (3.02115,-3.440) {\texttt{0000}};
\draw[trace] (6.04230,-3.440) -- (6.10230,-3.280) -- (10.94000,-3.280) -- (11.00000,-3.440) -- (10.94000,-3.600) -- (6.10230,-3.600) -- cycle;
\node at (8.52115,-3.440) {\texttt{FFF2}};
\node[anchor=east] at (-0.15,-4.040) {\texttt{sr}};
\draw[trace] (0.00000,-4.040) -- (0.06000,-3.880) -- (3.08286,-3.880) -- (3.14286,-4.040) -- (3.08286,-4.200) -- (0.06000,-4.200) -- cycle;
\node at (1.57143,-4.040) {\texttt{0000}};
\draw[trace] (3.14286,-4.040) -- (3.20286,-3.880) -- (10.94000,-3.880) -- (11.00000,-4.040) -- (10.94000,-4.200) -- (3.20286,-4.200) -- cycle;
\node at (7.07143,-4.040) {\texttt{B9B0}};
\node[anchor=east] at (-0.15,-4.640) {\texttt{dr}};
\draw[trace] (0.00000,-4.640) -- (0.06000,-4.480) -- (5.98230,-4.480) -- (6.04230,-4.640) -- (5.98230,-4.800) -- (0.06000,-4.800) -- cycle;
\node at (3.02115,-4.640) {\texttt{0000}};
\draw[trace] (6.04230,-4.640) -- (6.10230,-4.480) -- (10.94000,-4.480) -- (11.00000,-4.640) -- (10.94000,-4.800) -- (6.10230,-4.800) -- cycle;
\node at (8.52115,-4.640) {\texttt{B9B0}};
\node[anchor=east] at (-0.15,-5.240) {\texttt{valid}};
\draw[trace] (0.00000,-5.400) -- (6.04230,-5.400) -- (6.04230,-5.080) -- (7.10050,-5.080) -- (7.10050,-5.400) -- (11.00000,-5.400);
\draw[red!65!black,densely dashed] (6.04230,-5.550) -- (6.04230,.46);
\node[anchor=south,text=red!65!black] at (6.04230,.46) {$t_{\rm valid}$};
\end{tikzpicture}

\caption{本次姿态跨域仿真的首组快照波形（总线以十六进制表示）}
\label{fig:sim-cdc}
\end{figure}

100 次更新全部通过，无重复、丢失或字段混合。该测试的前提是源数据在目的端采样期间
保持稳定，且下一次更新留有足够间隔；数字 RTL 仿真不能模拟模拟电路的亚稳态概率，
相关跨域结构和时序约束仍需单独审查。
\sourcefile{\path{../Pgl50h/tests/tb_imu_protocol.sv}；
\path{../Pgl50h/tests/tb_attitude_cdc.sv}；
\path{chapter5/references/local_20261025/attitude_cdc.vcd}}

\section{像素预处理与帧级配置验证}
\subsection{Gamma 数值参考和视频节拍}
Gamma 测试使用数学公式独立计算表项期望值：
\begin{equation}
C_{\mathrm{ref}}=\operatorname{round}\!\left[255
\left(\frac{C_{\mathrm{in}}}{255}\right)^\gamma\right],
\qquad C_{\mathrm{in}}\in\{0,\ldots,255\}.
\end{equation}
七条有效曲线与一个保留恒等档共 2048 个表项全部比较，
不以“读取 DUT 的 LUT 再与 DUT 输出比较”作为数值参考。
像素计分板进一步检查三个颜色通道以及 DE、VS 的两级寄存器延迟，
包括增强关闭时的恒等像素行为与同样的节拍延迟。

场景激励包含暗场 16、亮场 224、明暗混合 32/224、中性灰 128、
各通道变化及 111/112 分类阈值附近的输入。测试检查分类滞回、
曲线逐档变化和帧内配置稳定性，主动在帧中改变启用请求，确认不会从半帧开始换曲线。
帧统计还检查末像素是否被计入、空帧和启动阶段是否被错误发布。
本次共核对 15365 个有效像素和 60 次统计，全部通过。

另有 10 月 22 日归档的 \path{tb_adaptive_gamma_1080p}，
在 148.5 MHz、60 Hz 完整光栅下检查 2073600 个像素不插入额外周期。
全白帧亮度和为
\begin{equation}
S_{\max}=255\times1920\times1080=528768000<2^{29},
\end{equation}
归档记录与此值精确一致。该结果用于补充整帧统计容量和连续节拍证据，
不是本次重新执行的第十二项测试。

\subsection{运行时掩码、SOF 原子应用和首像素对齐}
\path{tb_pc_control} 遍历三位功能掩码的全部 8 种组合，
bit0、bit1、bit2 分别控制镜头去畸变、姿态矫正和自适应 Gamma。
测试检查串口命令校验、未完成帧超时、busy 时拒收、状态查询与连续遥测，
重点区分“接收请求”和“输入 SOF 应用请求”两个时刻。
无 SOF 时 \path{requested_mask} 可以改变，
\path{applied_mask} 必须保持原值，并保留 pending 状态。

仅在简单测试平台中同时翻转掩码和 SOF，容易遗漏真实像素寄存器带来的相位差。
因此，本次还执行 \path{tb_pc_control_frame_alignment}，
串联输入寄存、控制跨域、Gamma 和 RGB565 转换，
在 SOF 与首个 DE 同时出现等条件下比较像素与帧元数据。
25 帧测试覆盖全部掩码与 pending 命令，检查首像素使用的 Gamma 开关、
几何控制和遥测帧身份一致，未发现一帧内混用新旧配置。

\begin{table}[htbp]\centering\small
\caption{帧级一致性检查及失败含义}\label{tab:sim-frame-atomic}
\begin{tabularx}{\textwidth}{>{\raggedright\arraybackslash}p{3.2cm}YY}
\toprule
检查事件 & 自动判定条件 & 可发现的问题 \\
\midrule
命令接收后尚无 SOF & 已应用配置不变，请求保持 pending & 串口 ACK 后立即切换半帧配置 \\
SOF 与首 DE 同时到达 & 首像素与该帧掩码使用相同配置 & 控制与像素流水线相差一拍 \\
帧内修改启用请求 & 当前完整帧的曲线和开关保持稳定 & 半帧亮度突变或错误统计归属 \\
Gamma 统计跨域 & 像素数、亮度和、曲线、掩码属于同一帧 & 分别同步多位字段导致快照撕裂 \\
\bottomrule
\end{tabularx}
\end{table}

\subsection{遥测与输出格式}
\path{tb_gamma_telemetry} 使用 \texttt{COLOR=1}，
独立计算彩色像素经 Gamma 后的 RGB 和近似亮度，
比较增强前后亮度和以及先前完成帧的曲线、掩码元数据，检查统计快照跨域时的一致性。
该检查防止出现“像素总数属于上一帧，而曲线号已属于下一帧”的情况。
统计位置在几何重采样和 RGB565 量化之前，其正确性不等于最终输出图像质量已获得验证。

输出格式测试穷举全部 65536 种 RGB565 输入，检查按位扩展至 RGB888 的结果及 DE 原样转发，
并分别检查默认 23 拍与参数化 2 拍光栅延迟，以及流水线占用期间的复位。
因此，格式测试的判据是像素位值与显示节拍一致，而非仅观察屏幕是否有图像。
\sourcefile{\path{../Pgl50h/tests/tb_adaptive_gamma_stream.sv}；
\path{../Pgl50h/tests/tb_pc_control_frame_alignment.sv}；
\path{../Pgl50h/tests/tb_gamma_telemetry.sv}；
\path{../Pgl50h/tests/tb_rgb888_raster.sv}；整帧 Gamma 归档见
\path{../Pgl50h/verification/pds_pc_control_write_ram_20261022/pc_control/03_tb_adaptive_gamma_1080p.log}}

\section{逆映射与分块渲染验证}
\subsection{坐标运算模块的参考点测试}
\path{tb_pose_inverse_project_q14} 采用默认除法实现，
先加载姿态并等待 \path{pose_ready}，再输入目的像素，
检查输出有效、源坐标和映射有效标志。
零姿态检查参考点 $(944,522)$，正负 yaw 约 $30^\circ$、pitch 约 $20^\circ$
以及组合姿态用于检查轴向、符号和投影方向；超限输入检查约 $35^\circ$ 的模块限幅，
并要求视场外坐标标记为无效，不能简单复制边界像素。

上述宽角度激励属于坐标模块测试，默认系统工作范围仍为 yaw/pitch 每轴约
$\pm15^\circ$。测试中参考点的允许误差为 2--5 像素，
因此通过结果只能说明这些参考点满足既定数值容差，
不能据此声称全视场亚像素精度、极低直线度误差或 MTF50 达标。

\subsection{分块遍历、缓存和逐像素计分板}
本次对 \path{tb_pose_renderer_tiled} 设置
\texttt{WIDTH=144, HEIGHT=40, TILE\_W=128, TILE\_H=32}，
保留其默认源图像 $320\times180$ 和测试内参 $f_x=f_y=125$、$(c_x,c_y)=(160,90)$。
这是截取左上方输出区域的功能测试，不是把完整标定图像缩放到 $144\times40$。
横向与纵向均产生一个完整块和一个裁剪块，四块分别为
$128\times32$、$16\times32$、$128\times8$、$16\times8$，
可直接覆盖右侧不足块宽和底部不足块高的边界。

源图像使用由 bank、横坐标和纵坐标共同决定的合成 RGB565 图样。
测试枚举所有目的像素，建立期望颜色数组和访问标志；每笔输出写突发到达时，
逐 beat、逐像素比较颜色，并拒绝越界、重复或遗漏输出。
源数据地址和输出地址还必须属于正确的帧 bank，突发长度必须处于 1--16 beats。
参考坐标由另一个相同的逆映射模块生成，故该计分板主要独立检查
调度、缓存寻址、采样和提交链路，不构成对投影公式的独立高精度复核。

渲染开始后，测试主动改变输入姿态、源 bank、输出 bank 和镜头使能请求，
要求当前任务继续使用启动时锁存的快照；读、写命令就绪分别周期性停顿，
响应延迟以 24 拍为基础附加变化，检查等待期间命令地址、长度和写载荷保持稳定。
只有全部输出写事务收到完成响应后才允许发布结果。

\begin{table}[htbp]\centering\small
\caption{本次 $144\times40$ 分块渲染测试结果}\label{tab:sim-tiled-results}
\begin{tabularx}{\textwidth}{rr*{5}{>{\centering\arraybackslash}X}}
\toprule
\multicolumn{2}{c}{姿态（Q2048）} & 耗时/拍 & 源读命令 & 源读 beats & 有效像素 & 黑色填充 \\
\cmidrule(lr){1-2}
yaw & pitch & & & & & \\
\midrule
0 & 0 & 12644 & 91 & 403 & 5760 & 0 \\
536 & 536 & 11552 & 23 & 129 & 866 & 4894 \\
536 & $-536$ & 12589 & 87 & 393 & 5760 & 0 \\
$-536$ & 536 & 10833 & 2 & 2 & 4 & 5756 \\
$-536$ & $-536$ & 12512 & 88 & 416 & 4279 & 1481 \\
\bottomrule
\end{tabularx}
\end{table}

表~\ref{tab:sim-tiled-results} 的 $536/2048$ rad 约为 $15^\circ$。
每组姿态均完成 5760 个目的像素、80 笔输出写事务；五组共比较 28800 个输出像素，
块边界、姿态快照和最后写入提交检查全部通过。
由于测试仅覆盖左上角区域，部分姿态下大多数坐标落到源视场之外；
黑色填充计数差异反映该测试几何条件，不能用来代表整幅输出的有效视场比例。
无效像素减少了源读取量，也使本表耗时不能直接外推到完整 1080p 帧率。

最后一次测试注入错误的读 LAST，确认模块报告错误后保持 busy、禁止 done，
防止错误帧被当作完整结果发布。
\sourcefile{\path{../Pgl50h/tests/tb_pose_inverse_project_q14.sv}；
\path{../Pgl50h/tests/tb_pose_renderer_tiled.sv}；
\path{chapter5/references/local_20261025/pose_renderer_tiled.log}}

\section{异步缓存与共享 DDR 协议验证}
\subsection{异步 FIFO 的顺序、反压和复位}
\path{tb_pose_async_fifo} 同时测试三种配置：256 bit、512 深度用于接近采集缓存的规模，
37 bit、8 深度以及 8 bit、2 深度用于快速触发满空边界和指针回绕。
各实例采用不同写读时钟，并插入随机停顿，由队列计分板核对读出顺序。
检查内容包含满时不写、空时不读、反压时输出保持、信用计数不高估可用空间，
以及 FIFO 非空时复位能够清除原有事务。

三组实例分别完成 12260/11488、15108/15096、3526/3522 次写/读，
所有检查通过。写读总数不同与测试中的非空复位及结束时残留有关，
不能直接把两者差值解释为丢失数据；无复位区间内的数据顺序由计分板逐项约束。

\subsection{DDR 仲裁与实际完成事件}
\path{tb_pose_ddr_arbiter} 设置 \texttt{STREAM=1}，
检查共享写载荷路径。采集写与渲染写、显示读与渲染读同时竞争时，
测试检查优先级、唯一事务归属和返回路由。行为模型独立停顿地址接收与数据传输，
插入最长 31 拍的延迟，避免使用固定计时器的错误完成逻辑侥幸通过。

工程所用 DDR 用户接口具有 AW/AR、WREADY、RVALID 和 LAST 等信号，
但没有完整 AXI4 的 WVALID、RREADY 与 B 响应通道；写侧 LAST 来自控制器。
因此，本章按该 IP 的接口约定判断有效数据拍，不套用完整 AXI4 的五通道握手规则。
每次有效写拍都比较载荷和 beat 序号，最后拍检查 LAST 与长度一致，
并要求完成信号只返回给该事务所属客户端。

图~\ref{fig:sim-ddr} 为本次仿真的三拍采集写片段，
\path{axi_awlen}=\texttt{2} 表示总计三个 beats。
755 ns 的上升沿接受地址，之后三个写数据拍分别在 815、935、1075 ns 被接受；
中间出现长时间 WREADY 停顿，\path{write_beat_index} 只能在真实数据拍后推进。
\path{cap_w_done} 于最后拍对应时刻产生，\path{protocol_error} 保持为零。
这组波形说明“地址已接受”与“最后数据已提交”是不同事件，
输出帧的发布不能由地址握手或估计延迟替代。

\begin{figure}[htbp]\centering
% Generated from actual VCD by references/make_waveforms.py.
\begin{tikzpicture}[x=1cm,y=1cm,font=\scriptsize]
\tikzset{trace/.style={draw=TemplateBlue,line width=0.65pt}}
\draw[gray!25] (0.00000,-4.350) -- (0.00000,0.75);
\node[anchor=south] at (0.00000,0.76) {740};
\draw[gray!25] (1.83333,-4.350) -- (1.83333,0.75);
\node[anchor=south] at (1.83333,0.76) {800};
\draw[gray!25] (3.66667,-4.350) -- (3.66667,0.75);
\node[anchor=south] at (3.66667,0.76) {860};
\draw[gray!25] (5.50000,-4.350) -- (5.50000,0.75);
\node[anchor=south] at (5.50000,0.76) {920};
\draw[gray!25] (7.33333,-4.350) -- (7.33333,0.75);
\node[anchor=south] at (7.33333,0.76) {980};
\draw[gray!25] (9.16667,-4.350) -- (9.16667,0.75);
\node[anchor=south] at (9.16667,0.76) {1040};
\draw[gray!25] (11.00000,-4.350) -- (11.00000,0.75);
\node[anchor=south] at (11.00000,0.76) {1100};
\node[anchor=west] at (11.15,0.87) {ns};
\node[anchor=east] at (-0.15,0.160) {\texttt{axi\_awvalid}};
\draw[trace] (0.00000,0.000) -- (0.15278,0.000) -- (0.15278,0.320) -- (0.45833,0.320) -- (0.45833,0.000) -- (10.54167,0.000) -- (10.54167,0.320) -- (10.84722,0.320) -- (10.84722,0.000) -- (11.00000,0.000);
\node[anchor=east] at (-0.15,-0.440) {\texttt{axi\_awready}};
\draw[trace] (0.00000,-0.280) -- (0.61111,-0.280) -- (0.61111,-0.600) -- (0.91667,-0.600) -- (0.91667,-0.280) -- (1.52778,-0.280) -- (1.52778,-0.600) -- (1.83333,-0.600) -- (1.83333,-0.280) -- (2.13889,-0.280) -- (2.13889,-0.600) -- (2.75000,-0.600) -- (2.75000,-0.280) -- (3.05556,-0.280) -- (3.05556,-0.600) -- (3.97222,-0.600) -- (3.97222,-0.280) -- (4.27778,-0.280) -- (4.27778,-0.600) -- (6.11111,-0.600) -- (6.11111,-0.280) -- (7.02778,-0.280) -- (7.02778,-0.600) -- (7.33333,-0.600) -- (7.33333,-0.280) -- (7.94444,-0.280) -- (7.94444,-0.600) -- (8.86111,-0.600) -- (8.86111,-0.280) -- (10.38889,-0.280) -- (10.38889,-0.600) -- (10.69444,-0.600) -- (10.69444,-0.280) -- (11.00000,-0.280);
\node[anchor=east] at (-0.15,-1.040) {\texttt{axi\_awlen}};
\draw[trace] (0.00000,-1.040) -- (0.05093,-0.880) -- (0.10185,-0.880) -- (0.15278,-1.040) -- (0.10185,-1.200) -- (0.05093,-1.200) -- cycle;
\draw[trace] (0.15278,-1.040) -- (0.21278,-0.880) -- (10.48167,-0.880) -- (10.54167,-1.040) -- (10.48167,-1.200) -- (0.21278,-1.200) -- cycle;
\node at (5.34722,-1.040) {\texttt{2}};
\draw[trace] (10.54167,-1.040) -- (10.60167,-0.880) -- (10.94000,-0.880) -- (11.00000,-1.040) -- (10.94000,-1.200) -- (10.60167,-1.200) -- cycle;
\node at (10.77083,-1.040) {\texttt{3}};
\node[anchor=east] at (-0.15,-1.640) {\texttt{axi\_wready}};
\draw[trace] (0.00000,-1.800) -- (2.13889,-1.800) -- (2.13889,-1.480) -- (2.44444,-1.480) -- (2.44444,-1.800) -- (5.80556,-1.800) -- (5.80556,-1.480) -- (6.11111,-1.480) -- (6.11111,-1.800) -- (10.08333,-1.800) -- (10.08333,-1.480) -- (10.38889,-1.480) -- (10.38889,-1.800) -- (11.00000,-1.800);
\node[anchor=east] at (-0.15,-2.240) {\texttt{axi\_wlast}};
\draw[trace] (0.00000,-2.400) -- (10.08333,-2.400) -- (10.08333,-2.080) -- (10.38889,-2.080) -- (10.38889,-2.400) -- (11.00000,-2.400);
\node[anchor=east] at (-0.15,-2.840) {\texttt{write\_beat\_index}};
\draw[trace] (0.00000,-2.840) -- (0.05093,-2.680) -- (0.10185,-2.680) -- (0.15278,-2.840) -- (0.10185,-3.000) -- (0.05093,-3.000) -- cycle;
\draw[trace] (0.15278,-2.840) -- (0.21278,-2.680) -- (2.23167,-2.680) -- (2.29167,-2.840) -- (2.23167,-3.000) -- (0.21278,-3.000) -- cycle;
\node at (1.22222,-2.840) {\texttt{0}};
\draw[trace] (2.29167,-2.840) -- (2.35167,-2.680) -- (5.89833,-2.680) -- (5.95833,-2.840) -- (5.89833,-3.000) -- (2.35167,-3.000) -- cycle;
\node at (4.12500,-2.840) {\texttt{1}};
\draw[trace] (5.95833,-2.840) -- (6.01833,-2.680) -- (10.48167,-2.680) -- (10.54167,-2.840) -- (10.48167,-3.000) -- (6.01833,-3.000) -- cycle;
\node at (8.25000,-2.840) {\texttt{2}};
\draw[trace] (10.54167,-2.840) -- (10.60167,-2.680) -- (10.94000,-2.680) -- (11.00000,-2.840) -- (10.94000,-3.000) -- (10.60167,-3.000) -- cycle;
\node at (10.77083,-2.840) {\texttt{0}};
\node[anchor=east] at (-0.15,-3.440) {\texttt{cap\_w\_done}};
\draw[trace] (0.00000,-3.280) -- (0.15278,-3.280) -- (0.15278,-3.600) -- (10.23611,-3.600) -- (10.23611,-3.280) -- (10.54167,-3.280) -- (10.54167,-3.600) -- (11.00000,-3.600);
\node[anchor=east] at (-0.15,-4.040) {\texttt{protocol\_error}};
\draw[trace] (0.00000,-4.200) -- (11.00000,-4.200);
\draw[red!65!black,densely dashed] (2.29167,-4.350) -- (2.29167,.46);
\node[anchor=south,text=red!65!black] at (2.29167,.46) {beat 0};
\draw[red!65!black,densely dashed] (5.95833,-4.350) -- (5.95833,.46);
\node[anchor=south,text=red!65!black] at (5.95833,.46) {beat 1};
\draw[red!65!black,densely dashed] (10.23611,-4.350) -- (10.23611,.46);
\node[anchor=south,text=red!65!black] at (10.23611,.46) {beat 2};
\end{tikzpicture}

\caption{三拍采集写事务在 DDR 反压下的完成时序（标线对应数据接受时钟沿）}
\label{fig:sim-ddr}
\end{figure}

本次正常事务共核对 160 笔写、760 个写 beats，以及 128 笔读、488 个读 beats。
此外，平台专门检查地址与首拍同周期的边界条件，并注入异常 LAST 和响应 ID，
要求错误被识别且完成事件不被错误路由。
上述协议验证通过后，仍需使用厂商模型检查真实控制器的突发边界与延迟分布。
\sourcefile{\path{../Pgl50h/tests/tb_pose_async_fifo.sv}；
\path{../Pgl50h/tests/tb_pose_ddr_arbiter.sv}；
\path{chapter5/references/local_20261025/pose_ddr_arbiter.vcd}}

\section{厂商模型与整帧并发归档复核}
\subsection{实际 DDR 控制器的突发边界测试}
10 月 10 日的 \path{tb_pose_vendor_ddr} 归档连接生成的
\path{DDR3_50H} 控制器、PHY 及两颗 Micron x16 DDR3 仿真模型。
测试覆盖 1--16 beats 的合法突发，并把起点放在 512/1024 个 32 bit 地址单位边界之前，
其中 1024 个地址单位对应 4 KiB。
为避免写入和读回共享相同的回绕错误，测试采用“突发写、单 beat 读回”和
“单 beat 写、突发读回”的非对称参考；另外检查 16 行、
行步长 960 个地址单位的源图像区域，等价于每行 1920 个 RGB565 像素。

归档共比较 472 个写 beats、608 个读 beats，数据和 LAST 位置检查通过。
100 MHz 用户时钟下，AR 接受到首个 RVALID 为 27--29 拍，
8-beat、16-beat 读到末拍分别为 34--41、42--49 拍，
AW 首次有效至最后 WREADY 的最大观测值为 245 拍。
模型中的跨行和成组刷新引入额外间隙，说明功能测试使用的 24 拍基础延迟
不能被当成真实 DDR 最坏访问延迟。

该归档未发现仿真 Error、Fatal 或 Micron 时序违规，存在 33 条与旧版生成 IP
端口、位宽和工具设置有关的 Warning。测试为短事务序列，
没有同时施加全速采集、渲染和显示负载，故不能单独证明整帧吞吐。
\sourcefile{\path{../Pgl50h/verification/pose_tiled_20261010/vendor_ddr_boundary_latency/README.md}；
同目录 \path{sim.log} 与 \path{source_identity.json}}

\subsection{完整 1080p 并发链路}
10 月 22 日归档的 \path{tb_fram_buf_pose_tiled} 在 ModelSim SE-64 10.6e 下运行，
输入为 $1920\times1080$ 完整光栅，共施加 28 个输入帧，启用预处理、运行时控制
与共享数据通道。DDR 行为模型设置基础延迟 32 拍、25\% 数据气泡，
并每 780 拍插入 35 拍刷新暂停；这与上一节厂商模型的实际短事务 trace 是两组不同条件。

\begin{table}[htbp]\centering\small
\caption{10 月 22 日整帧并发归档的功能观测}\label{tab:sim-full-archive}
\begin{tabularx}{\textwidth}{>{\raggedright\arraybackslash}p{4.2cm}Y}
\toprule
观测项 & 归档结果与解释 \\
\midrule
显示数据计分板 & 比较 51840000 个显示像素，检查帧身份及输出内容 \\
姿态与控制激励 & 11 组姿态激励及全部 8 种功能掩码；未穷举两者的笛卡尔积 \\
新结果发布 & 发布 13 帧；约 30.000030 fps，区别于 60 Hz 显示光栅 \\
最慢渲染与发布间隔 & 3314614 拍与 3333330 拍，按 100 MHz DDR 用户时钟计 \\
采集 FIFO 峰值 & 15/512，用于检查该场景下的缓冲压力 \\
日志状态 & 0 Error、4 Warning；警告涉及 Gamma 统计元数据未连接 \\
\bottomrule
\end{tabularx}
\end{table}

归档中的源帧编号会跳过部分输入帧，发布结果序列包含
1、2、5、7、\ldots、25；因此，通过判据是按系统策略发布完整且身份一致的新结果，
并非要求每个 60 Hz 输入帧都完成矫正。
该测试没有经真实 UART 注入全部控制流，也没有连接厂商 DDR PHY 模型。
后续源码已有修改，10 月 25 日复跑的小区域测试只能增加当前功能证据，
不能将上述 1080p 吞吐结果认定为当前全部源码的重测结论。
版本对应、帧率预算和性能余量进一步见第~\ref{chap:verification}~章。
\sourcefile{\path{../Pgl50h/verification/pds_pc_control_write_ram_20261022/simulation/result.log}；
同目录 \path{source_identity.json} 与 \path{../input_snapshot/}}

\section{复现方法与待补充验证}
在本报告根目录可执行
\begin{quote}\small\ttfamily
python3 chapter5/references/run\_simulations.py\\
python3 chapter5/references/make\_waveforms.py
\end{quote}
第一条命令运行表~\ref{tab:sim-matrix} 的 11 项测试并生成日志与清单，
第二条从 VCD 提取图~\ref{fig:sim-cdc}、图~\ref{fig:sim-ddr} 的真实信号跳变，
生成可随正文编译的 TikZ 波形。工程位置改变时，可通过
\texttt{--project} 指定路径；独立保留新结果可使用 \texttt{--output}。
波形生成脚本默认读取本章归档目录。

现有结果支持所选激励下的串口解析、跨域一致性、像素节拍、帧配置原子应用、
分块遍历和 DDR 完成语义。后续应在同一源码快照上重跑完整 1080p 并发压力测试，
补充复位与异常恢复、长时间多姿态切换及厂商 DDR 并发负载，
并将仿真快照与最终约束、实现报告和位流对应。
板级连续运行、端到端延迟和图卡清晰度仍须分别测量；
MTF50/SFR、直线度和边缘拉伸/断裂的定量判据与记录模板见第~\ref{chap:performance}~章。
